% answers_short.tex —— 2006 年·普通高等学校招生考试·山东卷文科 · 简短答案速查
% 由 tools/make_short_answers.py 自 answers.tex 的 \gksolhead/\gkanswer 提取，
% 并与 2006/2006_数学真题_山东文_答案要点.md 逐题比对后生成，勿手改（改源 MD 后重跑）。

\begin{gknotice}
\textbf{使用说明}：本节只给最终答案便于快速核对，完整推导见解析版。答案为编写者独立推导，\textbf{非官方标准答案}。
\end{gknotice}

\gkpart{选择题}
\begin{center}\begin{tabular}{ccccccccccccc}
\toprule
题号 & 1 & 2 & 3 & 4 & 5 & 6 & 7 & 8 & 9 & 10 & 11 & 12 \\
\midrule
答案 & D & C & A & D & B & B & A & C & A & D & A & B \\
\bottomrule
\end{tabular}\end{center}

\gkpart{填空题}
\begin{itemize}[leftmargin=2.6em]
  \item[\textbf{13.}] 150 人。
  \item[\textbf{14.}] $954/13$。
  \item[\textbf{15.}] 32。
  \item[\textbf{16.}] $\sqrt{21}/7$。
\end{itemize}

\gkpart{解答题}
\begin{itemize}[leftmargin=2.6em]
  \item[\textbf{17.}] (1) 若 $a>1$，递增区间为 $(-\infty,0)$、$(a-1,+\infty)$，递减区间为 $(0,a-1)$；若 $a=1$，在 $\mathbb R$ 上递增；(2) 若 $a>1$，极大值为 $1$，极小值为 $1-(a-1)^3$；若 $a=1$，无极值
  \item[\textbf{18.}] (1) $\pi/4$；(2) 2008
  \item[\textbf{19.}] (1) $9/14$；(2) $3/28$；(3) $4/7$
  \item[\textbf{20.}] (1) $2/\sqrt{15}$；(2) $45^\circ$；(3) $\lambda=2$
  \item[\textbf{21.}] (1) $x^2/2+y^2=1$；(2) $y=\sqrt{7/2}\,x+2$ 或 $y=-\sqrt{7/2}\,x+2$
  \item[\textbf{22.}] (1) $b_1=-3/4$，公比为 $1/2$；(2) $a_n=n-2+\dfrac3{2^n}$；(3) 存在，$\lambda=2$
\end{itemize}
